\chapter{Conclusion}
\label{ch:conclusion}

This study investigated the roles of temporal memory, risk-sensitive rewards, and turnover regularisation in deep reinforcement learning for daily portfolio allocation across ten major National Stock Exchange of India (NSE) equities over 2018--2025. Based on our four-fold, three-seed walk-forward experiments, we answer our core research questions:

\begin{itemize}
    \setlength{\itemsep}{1pt}
    \item \textbf{RQ1 (Temporal Architecture):} \textit{Does an LSTM recurrent encoder improve out-of-sample risk-adjusted returns and policy stability compared to an MLP baseline?}
    
    \textbf{Answer:} Yes. Under our setup, transitioning from the static MLP (M1) to the sequence-aware LSTM (M2) increased the pooled Sharpe ratio from $0.714$ to $0.915$ ($+0.201$, 95\% CI $[+0.156, +0.246]$), showing gains across all four folds ($4/4$). The recurrent representation cut average daily turnover from $0.186$ to $0.080$ and reduced cumulative transaction costs from $9.68\%$ to $4.27\%$ of capital, confirming that multi-day sequence features filter daily noise and prevent erratic reallocations.

    \item \textbf{RQ2 (Risk-Aware Reward Shaping):} \textit{Does transitioning from raw net return to the Differential Sharpe Ratio (DSR) improve Sharpe ratios and drawdowns?}
    
    \textbf{Answer:} The effect is mixed and regime-dependent. The DSR reward (M3) produced a flat incremental Sharpe delta ($+0.003$, 95\% CI $[-0.295, +0.222]$) relative to M2. While DSR performed strongly in trending regimes (Fold 3 Sharpe $1.328$ and Fold 4 Sharpe $1.164$), its recursive variance penalty slightly dampened pooled return ($26.93\%$ vs.\ $28.23\%$) and did not improve pooled maximum drawdown ($-19.95\%$ vs.\ $-18.52\%$).

    \item \textbf{RQ3 (Turnover Regularisation):} \textit{Does augmenting the reward with a quadratic penalty suppress turnover and transaction costs, and what is the trade-off?}
    
    \textbf{Answer:} Yes, with an explicit operational trade-off. Incorporating the quadratic weight-change penalty in Model M4 reduced average daily turnover from $0.079$ to $0.019$ (a 76\% reduction, improved in $12/12$ pairs) and lowered transaction fees to $0.97\%$ of capital. M4 also achieved the shallowest maximum drawdown among all evaluated models ($-16.04\%$, $+3.90\%$ over M3). However, this conservatism reduced pooled return from $26.93\%$ to $20.97\%$, demonstrating that penalising turnover restricts participation in rapid market rebounds.

    \item \textbf{RQ4 (Benchmark Competitiveness):} \textit{Can a continuous actor--critic DRL agent outperform a passively rebalanced equal-weight benchmark in a realistic Indian equity universe?}
    
    \textbf{Answer:} No. The passive Equal-Weight benchmark achieved superior pooled return ($36.48\%$) and Sharpe ratio ($1.13$). Our exploratory gap decomposition reveals that this is primarily driven by \textit{cash drag} ($-42.35\%$ in M4) resulting from learned cash reserves in a secular bull market under the zero-return cash assumption, alongside survivorship bias among established large-cap leaders.
\end{itemize}

In conclusion, deep reinforcement learning can autonomously acquire disciplined, low-turnover allocation strategies on continuous simplices. While active outperformance against fully-invested broad-market beta remains difficult due to cash drag and trading costs, our controlled ablation isolates the empirical value of sequence modeling, reward shaping, and turnover control for quantitative asset management.
